% GENERATED FILE. Edit canonical lessons, then run npm run generate:books.
% canonical-source-hash: fnv1a64:7f62b672195aeb20
% canonical-lessons: ES-C377-letra, ES-C377-frase, ES-C377-idioma, ES-C377-historia, ES-C377-baile

\chapter{Five Things You Can Learn}
\label{ch:la-letra-y-la-frase}
\hlchaptermodality{377}

\begin{chapteropening}
\textbf{By the end of this chapter:} \emph{I can say la letra, la frase, el idioma, la historia and el baile, and name five things a person can learn.}

Complete ES-C377-baile: name five things a person can learn, and say one thing about each.
\end{chapteropening}

\section[la letra]{la letra — the word for writing, and nobody knows where it came from}

\label{lesson:ES-C377-letra}



\begin{quote}
Say \emph{the chocolate}, then \emph{the pepper}. (\emph{El chocolate}; \emph{El pimiento}.)
\end{quote}

\subsection*{What to know first}
\begin{itemize}
\raggedright
  \item la palabra — what letters build.
  \item la carta — the other thing English calls a letter.
\end{itemize}

\begin{sounds}
\begin{itemize}
\raggedright
  \item \emph{LE-tra}, two beats with the weight on the first. The \emph{tr} runs together and the \emph{r} is a single tap. $\to$ pronunciation reference
\end{itemize}
\end{sounds}

\begin{cousinweb}
\textbf{la letra} — \textquotedblleft{}letter.\textquotedblright{}

Latin \textbf{littera} was a letter of the alphabet, and English took it three times:

\begin{itemize}
\raggedright
  \item \textbf{letter}, through French, worn down;
  \item \textbf{literal}, straight from Latin, meaning by the letter;
  \item \textbf{literature}, the whole of what letters have been used for.
\end{itemize}

Watch out for one thing in English. \emph{Letter} covers two different objects: a mark on a page and something you post. Spanish keeps them apart. \emph{La letra} is the mark. The thing you post is \textbf{la carta}, which you already have.

And here is the part that ought to be embarrassing and is only honest. Nobody knows where Latin got \emph{littera}.

Two attempts exist. One builds it on \textbf{linō}, \textquotedblleft{}to smear\textquotedblright{} — writing as smearing — and it fails on the shape of the word rather than the sense. The other has Latin borrowing it from Greek \textbf{diphthéra}, a prepared hide people wrote on, by way of Etruscan. That one is still live and still unproven.

The word for writing is one of the words we cannot trace. Whatever else literacy did, it did not preserve its own name.
\end{cousinweb}

\begin{grammarlens}[title={What you've built}]
You can now name a single character: \emph{una letra}.
\end{grammarlens}

\subsection*{Your turn}
Say these aloud:

\begin{itemize}
\raggedright
  \item \emph{la letra}
  \item \emph{una letra}
  \item \emph{la primera letra}
\end{itemize}

\begin{tcolorbox}[breakable,colback=teal!4,colframe=teal!35!black,title={Before you move on}]
Which Spanish word means a letter you post? (\emph{La carta} — not \emph{la letra}.) And does anyone know where Latin got \emph{littera}? (No; both proposals are unproven.)
\end{tcolorbox}
\section[la frase]{la frase — a way of speaking}

\label{lesson:ES-C377-frase}



\begin{quote}
Say \emph{the pepper}, then \emph{the letter}. (\emph{El pimiento}; \emph{La letra}.)
\end{quote}

\subsection*{What to know first}
\begin{itemize}
\raggedright
  \item la palabra — the smaller unit.
  \item hablar — what you do with one.
\end{itemize}

\begin{sounds}
\begin{itemize}
\raggedright
  \item \emph{FRA-se}, two beats with the weight on the first. The \emph{fr} runs together and the \emph{r} is a single tap. $\to$ pronunciation reference
\end{itemize}
\end{sounds}

\begin{cousinweb}
\textbf{la frase} — \textquotedblleft{}phrase.\textquotedblright{}

Greek \textbf{phrásis} meant an expression, a manner of speaking. Behind it is the verb \textbf{phrázein}, \textquotedblleft{}to point out, to declare.\textquotedblright{} The word came into Latin as \emph{phrasis} and into Spanish from there.

English has it twice, and the second one repays unpacking. A \textbf{paraphrase} is \emph{para-} \textquotedblleft{}beside\textquotedblright{} plus \emph{phrasis}: a saying-alongside. When you paraphrase, you set a second version of the words next to the first. The word describes the action exactly.

One useful difference from English. Spanish \emph{frase} is wider than English \emph{phrase} — it covers a whole sentence, not only a fragment of one. \emph{Una frase} is something you can say and stop.
\end{cousinweb}

\begin{grammarlens}[title={What you've built}]
You can now name a whole thing you say: \emph{una frase}.
\end{grammarlens}

\subsection*{Your turn}
Say these aloud:

\begin{itemize}
\raggedright
  \item \emph{la frase}
  \item \emph{una frase}
  \item \emph{una frase en español}
\end{itemize}

\begin{tcolorbox}[breakable,colback=teal!4,colframe=teal!35!black,title={Before you move on}]
What did Greek \emph{phrásis} mean? (An expression, a way of speaking.) And what does \emph{para-} add in \emph{paraphrase}? (Beside — a saying-alongside.)
\end{tcolorbox}
\section[el idioma]{el idioma — what belongs to you}

\label{lesson:ES-C377-idioma}



\begin{quote}
Say \emph{the letter}, then \emph{the phrase}. (\emph{La letra}; \emph{La frase}.)
\end{quote}

\subsection*{What to know first}
\begin{itemize}
\raggedright
  \item español — the one you are learning.
  \item la palabra — what one is made of.
\end{itemize}

\begin{sounds}
\begin{itemize}
\raggedright
  \item \emph{i-DIO-ma}, three beats with the weight on the middle one. The \emph{io} is one glide. The word ends in \emph{-a} and is masculine — \emph{el idioma}. $\to$ pronunciation reference
\end{itemize}
\end{sounds}

\begin{cousinweb}
\textbf{el idioma} — \textquotedblleft{}language.\textquotedblright{}

Greek \textbf{ídios} meant \textquotedblleft{}one's own, private, particular.\textquotedblright{} From it Greek built \textbf{idíōma}, a peculiarity of style — the way of speaking that belongs to one person or one place. Late Latin took it, and Spanish took it from there.

Three English words come off that same root, and the roads are worth walking separately:

\begin{itemize}
\raggedright
  \item an \textbf{idiom} is a turn of phrase belonging to one language and refusing to move into another;
  \item an \textbf{idiosyncrasy} is a temperament belonging to one person;
  \item an \textbf{idiot} is, in Greek, an \textbf{idiōtēs} — a \textbf{private} person. Someone who held no public office and took no part in civic life. From there it slid to \textquotedblleft{}layman,\textquotedblright{} then to \textquotedblleft{}non-expert,\textquotedblright{} and only much later to the insult.
\end{itemize}

So the word for a language and the word for a fool start in the same place: what is yours alone. One kept the pride in that and the other picked up the contempt.

A grammatical note you will need. \emph{El idioma} ends in \emph{-a} and is masculine. That \emph{-ma} ending is Greek, and Spanish keeps those nouns masculine as a group.
\end{cousinweb}

\begin{grammarlens}[title={What you've built}]
You can now name what you are learning: \emph{un idioma}.
\end{grammarlens}

\subsection*{Your turn}
Say these aloud:

\begin{itemize}
\raggedright
  \item \emph{el idioma}
  \item \emph{un idioma}
  \item \emph{el idioma español}
\end{itemize}

\begin{tcolorbox}[breakable,colback=teal!4,colframe=teal!35!black,title={Before you move on}]
What did Greek \emph{ídios} mean? (One's own, private.) And what was an \emph{idiōtēs} in Greek? (A private person, one with no public role.)
\end{tcolorbox}
\section[la historia]{la historia — one Greek word that became two English ones}

\label{lesson:ES-C377-historia}



\begin{quote}
Say \emph{the phrase}, then \emph{the language}. (\emph{La frase}; \emph{El idioma}.)
\end{quote}

\subsection*{What to know first}
\begin{itemize}
\raggedright
  \item la carta — one place you might write one down.
  \item la palabra — what it is built from.
\end{itemize}

\begin{sounds}
\begin{itemize}
\raggedright
  \item \emph{his-TO-ria}, three beats with the weight on the middle one. The \emph{h} is silent, so the word opens on the \emph{i}. $\to$ pronunciation reference
\end{itemize}
\end{sounds}

\begin{cousinweb}
\textbf{la historia} — \textquotedblleft{}story.\textquotedblright{}

Greek \textbf{historía} did not mean the past. It meant \textbf{inquiry} — and then the knowledge you come away with after inquiring. A \emph{hístōr} was someone who knows because they went and found out.

The narrowing to \textquotedblleft{}past events\textquotedblright{} happened later, and it is worth keeping in mind, because it turns \emph{history} from a subject into a method.

English then took the word \textbf{twice}:

\begin{itemize}
\raggedright
  \item \textbf{history}, borrowed formally from Latin \emph{historia};
  \item \textbf{story}, which came through Anglo-French \emph{estorie} and lost its opening syllable on the way.
\end{itemize}

For a long time English used them interchangeably. The split between a \emph{history} that happened and a \emph{story} that might not is only about five hundred years old.

Spanish never made that split. \textbf{Una historia} is both — the account of what happened and the tale you tell a child — and which one you mean comes from the situation rather than the word.
\end{cousinweb}

\begin{grammarlens}[title={What you've built}]
You can now name what you tell: \emph{una historia}.
\end{grammarlens}

\subsection*{Your turn}
Say these aloud:

\begin{itemize}
\raggedright
  \item \emph{la historia}
  \item \emph{una historia}
  \item \emph{una historia larga}
\end{itemize}

\begin{tcolorbox}[breakable,colback=teal!4,colframe=teal!35!black,title={Before you move on}]
What did Greek \emph{historía} originally mean? (Inquiry, and the knowledge you get from it.) And which English word is the same one worn down through French? (\emph{Story}.)
\end{tcolorbox}
\section[el baile]{el baile — two English balls, one of them unrelated}

\label{lesson:ES-C377-baile}



\begin{quote}
Say \emph{the language}, then \emph{the story}. (\emph{El idioma}; \emph{La historia}.)
\end{quote}

\subsection*{What to know first}
\begin{itemize}
\raggedright
  \item la fiesta — where one happens.
  \item la risa — what usually comes with it.
\end{itemize}

\begin{sounds}
\begin{itemize}
\raggedright
  \item \emph{BAI-le}, two beats with the weight on the first. The \emph{ai} is one glide, close to English \emph{eye}. $\to$ pronunciation reference
\end{itemize}
\end{sounds}

\begin{cousinweb}
\textbf{el baile} — \textquotedblleft{}dance.\textquotedblright{}

Late Latin \textbf{ballāre} meant to dance. What Spanish did with it is not the straightforward thing.

An inherited \emph{ballāre} would have come out as \emph{\textbf{ballar}}. Spanish says \textbf{bailar}, and Corominas reads that as a \textbf{borrowing} — from Occitan \textbf{balar}, taken up around 1270. The stray \emph{-i-} has never been fully explained.

Go one step further back and the Academy's dictionary hedges twice. It offers Greek \textbf{pállein}, \textquotedblleft{}to shake,\textquotedblright{} and marks the link \emph{perhaps} both times. So the chain is: certain Spanish word, probable Occitan route, secure Late Latin verb, speculative Greek origin — four links of steadily decreasing confidence.

English inherited the middle of that chain and shows it off:

\begin{itemize}
\raggedright
  \item a \textbf{ball}, the event you dress up for;
  \item a \textbf{ballet}, through Italian;
  \item a \textbf{ballad}, which began as a song to dance to.
\end{itemize}

And then there is the other \textbf{ball}, the round one you throw. That is a Germanic word with no connection whatsoever. Two identical English words, two separate ancestries, and the one that means a dance is the one related to \emph{baile}.
\end{cousinweb}

\begin{grammarlens}[title={What you've built}]
You can now name what happens at a party: \emph{un baile}.
\end{grammarlens}

\subsection*{Your turn}
Say these aloud:

\begin{itemize}
\raggedright
  \item \emph{el baile}
  \item \emph{un baile}
  \item \emph{un baile largo}
\end{itemize}

\begin{tcolorbox}[breakable,colback=teal!4,colframe=teal!35!black,title={Before you move on}]
Did Spanish inherit \emph{bailar} from Latin or borrow it? (Borrow it — from Occitan.) And which English \emph{ball} is related to it? (The dance, not the round thing.)
\end{tcolorbox}
